\documentclass[10pt,a4paper]{amsart}
\usepackage[margin=19mm]{geometry}
\usepackage{amsmath,amssymb,mathtools}
\usepackage[scr=rsfs,cal=euler]{mathalfa}
\usepackage{xfrac}
\usepackage[unicode,hidelinks,bookmarks=false]{hyperref}
\newcommand{\ie}{i.e.\ }
\newcommand{\cf}{cf.\ }
\newcommand{\resp}{respectively\ }
\newcommand{\Subsection}{\subsection}
\providecommand{\nobreakdash}{\nobreak\hbox{-}}
\setlength{\emergencystretch}{2em}
\multlinegap0pt
\usepackage{bm}
\usepackage{tikz}
\usepackage{mathtools}
\usepackage[shortlabels]{enumitem}
\usepackage{amssymb}
\newtheorem{proposition}{Proposition}
\newcommand{\R}{\mathbb{R}}
\newcommand{\mB}{\mathcal{B}}
\newcommand{\wF}{\widehat{F}}
\title{On the computation of fold and cusp bifurcations in mass action networks via dual equations}
\author{Murad Banaji}
\date{Source: arXiv:2609.07265v1 (CC BY 4.0)}
\begin{document}
\maketitle

\noindent\emph{Source and licence.} On the computation of fold and cusp bifurcations in mass action networks via dual equations. Version arXiv:2609.07265v1 (CC BY 4.0), distributed by arXiv under the Creative Commons Attribution 4.0 International licence, http://creativecommons.org/licenses/by/4.0/, as recorded in the arXiv licence metadata for this exact version and shown on the abstract page https://arxiv.org/abs/2609.07265v1. Full licence notice: \texttt{licenses/arxiv-2609.07265-CC-BY-4.0.txt}.\par\smallskip

\noindent\textit{Editorial scope.} Counted unit: the article's proposition proposition (source0.tex lines 280--287). The article's complete original proof of it is delivered with it (source0.tex lines 288--296, one \texttt{proof} environment of 9 lines). No mathematics is written, repaired or re-derived: the statement and the proof are the article's own lines, in the article's own notation, and no retained line is substituted or re-flowed.  No further source-local context is needed: every cross-reference of every retained line resolves inside the two delivered spans.  Nothing was omitted from any retained span, so the package carries no omission record.  The retained lines cite 1 works, whose entries are transcribed verbatim from the article's own bibliography source in the e-print and delivered with short editorial labels recorded in the item's reference mapping.  No retained line contains a figure environment, an image-inclusion command or drawing code, so no image is delivered and the package declares no asset dependency.  Editorial formatting only, in addition: an AMS \texttt{amsart} preamble that restates the article's own declarations and macros used by the retained lines, namely the environments \texttt{proposition} on separate counters, together with the standard packages those lines need.  The retained environments print in this document as Proposition 1 in order of appearance, whereas the article prints the same items with the numbering of its own full document.  The complete native source of the article as distributed in the arXiv e-print for this exact version is bundled unchanged as \texttt{sources/209615/source0.tex}.
\begin{proposition}[Reparameterisation]
\label[proposition]{prop:reparam}
Suppose that the system $\dot y = f(y,\beta)$ has some local bifurcation of equilibria $\mB$ at $(y_0, \beta_0)$. Let $k \leq m$ and $\Theta \colon V \to \R^k$ be any smooth map such that $\Theta_\beta(\beta_0)$ is surjective, and suppose that $f(y,\beta) \equiv \hat{f}(y, \Theta(\beta))$, where this equation defines $\hat{f}\colon U \times \Theta(V) \to \R^n$. Then the system $\dot y = \hat{f}(y,\theta)$ has the bifurcation $\mB$ at $(y_0, \theta_0)$, where $\theta_0:=\Theta(\beta_0)$ and the following are equivalent:
\begin{enumerate}
\item[(a)] The bifurcation in $\dot y = f(y,\beta)$ at $(y_0, \beta_0)$ is unfolded transversally by $\beta$;
\item[(b)] The bifurcation in $\dot y = \hat{f}(y,\theta)$ at $(y_0, \theta_0)$ is unfolded transversally by $\theta$.
\end{enumerate}
\end{proposition}

\begin{proof}
The claim that $\dot y = \hat{f}(y,\theta)$ has the bifurcation $\mB$ at $(y_0, \theta_0)$ follows trivially from our general assumptions about bifurcations, as the Taylor coefficients of $f$ at $(y_0, \beta_0)$ and $\hat{f}$ at  $(y_0, \theta_0)$ w.r.t. $y$ are identical to all orders. As noted in \cite{banaji:boros:hofbauer:2025}, associated with the bifurcation $\mB$ in $\dot y = f(y,\beta)$ is a ``bifurcation function'', say $F(y, \beta)$, such that in a neighbourhood of $(y_0, \beta_0)$, the bifurcation occurs if and only if $F(y,\beta)=0$. Further, the bifurcation is unfolded transversally by the parameters $\beta$ if and only if (i) $F_y(y_0, \beta_0)$ is injective and (ii) $\partial_{(y,\beta)} F(y_0, \beta_0)$ is surjective. Note that the bifurcation function $\wF(y,\theta)$ for the reparameterised system $\dot y = \hat{f}(y,\theta)$ is naturally defined via $F(y,\beta) = \wF(y,\Theta(\beta))$.

As $F_y(y_0, \beta_0) = \wF_y(y_0, \theta_0)$, the first condition for transversality holds in one case if and only if it holds in the other. Moreover, by the chain rule,
\[
[F_y\,|\, F_\beta] = [\wF_y\,|\,\wF_\theta]M \quad \mbox{where} \quad M=\left(\begin{array}{cc}I&0\\0& \Theta_\beta(\beta_0)\end{array}\right)\,,
\]
hence $\mathrm{rank}\,[F_y\,|\, F_\beta] \leq \mathrm{rank}[\wF_y\,|\,\wF_\theta]$. On the other hand, $M$ is a surjective matrix and hence has a right inverse, say $M^*$, giving $[F_y\,|\, F_\beta]M^* = [\wF_y\,|\,\wF_\theta]$; consequently, $\mathrm{rank}\,[F_y\,|\, F_\beta] \geq \mathrm{rank}[\wF_y\,|\,\wF_\theta]$. Thus $\mathrm{rank}\,[F_y\,|\, F_\beta] = \mathrm{rank}[\wF_y\,|\,\wF_\theta]$, and the second condition for transversality holds in one case if and only if it holds in the other. 
\end{proof}

\begin{thebibliography}{99}
\bibitem[banaji]{banaji:boros:hofbauer:2025} M.~Banaji, B.~Boros, and J.~Hofbauer. \newblock The inheritance of local bifurcations in mass action networks. \newblock {\em J. Nonlinear Sci.}, 35(72), 2025.
\end{thebibliography}
\end{document}
